\documentclass[12pt]{article}
\usepackage{amsmath,amssymb}

\begin{document}
\section*{{2022 AIME II Problems/Problem 12}}
Source: AoPS Wiki

\subsection*{Solution}
Denote $P = \left( x , y \right)$ . Because $\frac{x^2}{a^2}+\frac{y^2}{a^2-16} = 1$ , $P$ is on an ellipse whose center is $\left( 0 , 0 \right)$ and foci are $\left( - 4 , 0 \right)$ and $\left( 4 , 0 \right)$ . Hence, the sum of distance from $P$ to $\left( - 4 , 0 \right)$ and $\left( 4 , 0 \right)$ is equal to twice the major axis of this ellipse, $2a$ . Because $\frac{(x-20)^2}{b^2-1}+\frac{(y-11)^2}{b^2} = 1$ , $P$ is on an ellipse whose center is $\left( 20 , 11 \right)$ and foci are $\left( 20 , 10 \right)$ and $\left( 20 , 12 \right)$ . Hence, the sum of distance from $P$ to $\left( 20 , 10 \right)$ and $\left( 20 , 12 \right)$ is equal to twice the major axis of this ellipse, $2b$ . Therefore, $2a + 2b$ is the sum of the distance from $P$ to four foci of these two ellipses. To make this minimized, $P$ is the intersection point of the line that passes through $\left( - 4 , 0 \right)$ and $\left( 20 , 10 \right)$ , and the line that passes through $\left( 4 , 0 \right)$ and $\left( 20 , 12 \right)$ . The distance between $\left( - 4 , 0 \right)$ and $\left( 20 , 10 \right)$ is $\sqrt{\left( 20 + 4 \right)^2 + \left( 10 - 0 \right)^2} = 26$ . The distance between $\left( 4 , 0 \right)$ and $\left( 20 , 12 \right)$ is $\sqrt{\left( 20 - 4 \right)^2 + \left( 12 - 0 \right)^2} = 20$ . Hence, $2 a + 2 b = 26 + 20 = 46$ . Therefore, $a + b = \boxed{\textbf{(023) }}.$ ~Steven Chen (www.professorchenedu.com)

\end{document}
